% GENERATED FILE. Edit canonical lessons, then run npm run generate:books.
% canonical-source-hash: fnv1a64:7401bcd7df4201e9
% canonical-lessons: LA-C418-aes-alienum, LA-C418-debitor, LA-C418-creditor, LA-C418-fenus, LA-C418-mutuus

\chapter{Debts}
\label{ch:la-debts}
\hlchaptermodality{418}

\begin{chapteropening}
\textbf{By the end of this chapter:} \emph{I can talk about debt, a debtor and a creditor, interest, and money borrowed or lent.}

Complete the last lesson of chapter 418: I can talk about debt, a debtor and a creditor, interest, and money borrowed or lent.
\end{chapteropening}

\section[aes aliēnum, aeris aliēnī]{aes aliēnum, aeris aliēnī — debt}

\label{lesson:LA-C418-aes-alienum}



\begin{quote}
Before the new one: say the Latin for ownership, then the Latin for formal purchase.
\end{quote}

\subsection*{You'll want to know: aes aliēnum, aeris aliēnī}
\textbf{aes aliēnum, aeris aliēnī} — \textquotedblleft{}debt\textquotedblright{}.

\begin{quote}
Aere aliēnō premitur.
\end{quote}

\begin{quote}
\emph{He is weighed down by debt.}
\end{quote}

\emph{aes aliēnum, aeris aliēnī} n.: other people's bronze, money that belongs to someone else. \emph{In aere aliēnō esse}, to be in debt.

\begin{cousinweb}
\emph{aes aliēnum} is built from \emph{aes} \textquotedblleft{}bronze, money\textquotedblright{} and \emph{aliēnus} \textquotedblleft{}belonging to another\textquotedblright{} (Lewis \& Short): money that is someone else's.
\end{cousinweb}

\begin{grammarlens}[title={What you've built}]
One more word for telling, explaining and describing at B1.
\end{grammarlens}

\subsection*{Your turn}
\begin{itemize}
\raggedright
  \item \emph{Say it:} \emph{aes aliēnum, aeris aliēnī}
  \item \emph{Say it:} \emph{aes aliēnum, aeris aliēnī}, once more
  \item \emph{Say it:} a sentence of your own with \emph{aes aliēnum, aeris aliēnī}
  \item \emph{Recall:} say the Latin for ownership, then the Latin for formal purchase, then say \emph{aes aliēnum, aeris aliēnī} again
\end{itemize}

\begin{tcolorbox}[breakable,colback=teal!4,colframe=teal!35!black,title={Before you move on}]
What does aes aliēnum, aeris aliēnī mean? (\textquotedblleft{}Debt\textquotedblright{}.) Say it once more.
\end{tcolorbox}
\section[dēbitor, dēbitōris]{dēbitor, dēbitōris — a debtor}

\label{lesson:LA-C418-debitor}



\begin{quote}
Before the new one: say the Latin for formal purchase, then the Latin for debt.
\end{quote}

\subsection*{You'll want to know: dēbitor, dēbitōris}
\textbf{dēbitor, dēbitōris} — \textquotedblleft{}a debtor\textquotedblright{}.

\begin{quote}
Dēbitor pecūniam reddere nōn potuit.
\end{quote}

\begin{quote}
\emph{The debtor could not pay the money back.}
\end{quote}

\emph{dēbitor, dēbitōris} m., built on \emph{dēbeō}, to owe. English \emph{debtor}.

\begin{cousinweb}
\emph{dēbitor} is built on \emph{dēbitum} \textquotedblleft{}a debt\textquotedblright{} (EtymDB), from \emph{dēbeō} \textquotedblleft{}to owe\textquotedblright{} (Lewis \& Short). English \emph{debtor} comes from it through French.
\end{cousinweb}

\begin{grammarlens}[title={What you've built}]
One more word for telling, explaining and describing at B1.
\end{grammarlens}

\subsection*{Your turn}
\begin{itemize}
\raggedright
  \item \emph{Say it:} \emph{dēbitor, dēbitōris}
  \item \emph{Say it:} \emph{dēbitor, dēbitōris}, once more
  \item \emph{Say it:} a sentence of your own with \emph{dēbitor, dēbitōris}
  \item \emph{Recall:} say the Latin for formal purchase, then the Latin for debt, then say \emph{dēbitor, dēbitōris} again
\end{itemize}

\begin{tcolorbox}[breakable,colback=teal!4,colframe=teal!35!black,title={Before you move on}]
What does dēbitor, dēbitōris mean? (\textquotedblleft{}A debtor\textquotedblright{}.) Say it once more.
\end{tcolorbox}
\section[crēditor, crēditōris]{crēditor, crēditōris — a creditor, a lender}

\label{lesson:LA-C418-creditor}



\begin{quote}
Before the new one: say the Latin for debt, then the Latin for a debtor.
\end{quote}

\subsection*{You'll want to know: crēditor, crēditōris}
\textbf{crēditor, crēditōris} — \textquotedblleft{}a creditor, a lender\textquotedblright{}.

\begin{quote}
Crēditor pecūniam suam exspectat.
\end{quote}

\begin{quote}
\emph{The creditor is waiting for his money.}
\end{quote}

\emph{crēditor, crēditōris} m., built on \emph{crēdō}, to trust, to lend. English \emph{creditor}.

\begin{cousinweb}
\emph{crēditor} is built on \emph{crēdō} \textquotedblleft{}to trust, to lend\textquotedblright{} (EtymDB; Lewis \& Short), from its participle \emph{crēditus}. English \emph{creditor} comes from it.
\end{cousinweb}

\begin{grammarlens}[title={What you've built}]
One more word for telling, explaining and describing at B1.
\end{grammarlens}

\subsection*{Your turn}
\begin{itemize}
\raggedright
  \item \emph{Say it:} \emph{crēditor, crēditōris}
  \item \emph{Say it:} \emph{crēditor, crēditōris}, once more
  \item \emph{Say it:} a sentence of your own with \emph{crēditor, crēditōris}
  \item \emph{Recall:} say the Latin for debt, then the Latin for a debtor, then say \emph{crēditor, crēditōris} again
\end{itemize}

\begin{tcolorbox}[breakable,colback=teal!4,colframe=teal!35!black,title={Before you move on}]
What does crēditor, crēditōris mean? (\textquotedblleft{}A creditor, a lender\textquotedblright{}.) Say it once more.
\end{tcolorbox}
\section[fēnus, fēnoris]{fēnus, fēnoris — interest (on money)}

\label{lesson:LA-C418-fenus}



\begin{quote}
Before the new one: say the Latin for a debtor, then the Latin for a creditor.
\end{quote}

\subsection*{You'll want to know: fēnus, fēnoris}
\textbf{fēnus, fēnoris} — \textquotedblleft{}interest (on money)\textquotedblright{}.

\begin{quote}
Pecūniam fēnore dabat.
\end{quote}

\begin{quote}
\emph{He lent money at interest.}
\end{quote}

\emph{fēnus, fēnoris} n., also spelled \emph{faenus}: the yield of money, as a field yields a crop.

\begin{cousinweb}
\emph{fēnus} is usually traced to a root meaning \textquotedblleft{}to bring forth, to yield\textquotedblright{}, the root Lewis \& Short also see in \emph{fēmina}: what money produces. The account is not certain.
\end{cousinweb}

\begin{grammarlens}[title={What you've built}]
One more word for telling, explaining and describing at B1.
\end{grammarlens}

\subsection*{Your turn}
\begin{itemize}
\raggedright
  \item \emph{Say it:} \emph{fēnus, fēnoris}
  \item \emph{Say it:} \emph{fēnus, fēnoris}, once more
  \item \emph{Say it:} a sentence of your own with \emph{fēnus, fēnoris}
  \item \emph{Recall:} say the Latin for a debtor, then the Latin for a creditor, then say \emph{fēnus, fēnoris} again
\end{itemize}

\begin{tcolorbox}[breakable,colback=teal!4,colframe=teal!35!black,title={Before you move on}]
What does fēnus, fēnoris mean? (\textquotedblleft{}Interest (on money)\textquotedblright{}.) Say it once more.
\end{tcolorbox}
\section[mūtuus, mūtua, mūtuum]{mūtuus, mūtua, mūtuum — borrowed, lent; mutual}

\label{lesson:LA-C418-mutuus}



\begin{quote}
Before the new one: say the Latin for a creditor, then the Latin for interest.
\end{quote}

\subsection*{You'll want to know: mūtuus, mūtua, mūtuum}
\textbf{mūtuus, mūtua, mūtuum} — \textquotedblleft{}borrowed, lent; mutual\textquotedblright{}.

\begin{quote}
Pecūniam mūtuam ab amīcō sūmpsit.
\end{quote}

\begin{quote}
\emph{He borrowed money from a friend.}
\end{quote}

\emph{mūtuus, mūtua, mūtuum}: \emph{pecūniam mūtuam dare}, to lend; \emph{sūmere}, to borrow. Also mutual: \emph{amor mūtuus}. English \emph{mutual}.

\begin{cousinweb}
\emph{mūtuus} is built on \emph{mūtō} \textquotedblleft{}to change, to exchange\textquotedblright{} (EtymDB; Lewis \& Short): given in exchange. English \emph{mutual} comes from it.
\end{cousinweb}

\begin{grammarlens}[title={What you've built}]
One more word for telling, explaining and describing at B1.
\end{grammarlens}

\subsection*{Your turn}
\begin{itemize}
\raggedright
  \item \emph{Say it:} \emph{mūtuus, mūtua, mūtuum}
  \item \emph{Say it:} \emph{mūtuus, mūtua, mūtuum}, once more
  \item \emph{Say it:} a sentence of your own with \emph{mūtuus, mūtua, mūtuum}
  \item \emph{Recall:} say the Latin for a creditor, then the Latin for interest, then say \emph{mūtuus, mūtua, mūtuum} again
\end{itemize}

\begin{tcolorbox}[breakable,colback=teal!4,colframe=teal!35!black,title={Before you move on}]
What does mūtuus, mūtua, mūtuum mean? (\textquotedblleft{}Borrowed, lent; mutual\textquotedblright{}.) Say it once more.
\end{tcolorbox}
